\section{Entropy bounds in a Hermitian background}

For a Hermitian background and reference zero, the dual functional need not equal entropy plus a constant. Its finite domain is nevertheless the finite-entropy domain.

\begin{proposition}\label{prop:hermitian-dual-entropy-bounds}
Let $\theta=\omega$ be Hermitian, let $\varphi=0$, and normalize $\mathrm{d}V(X)=1$. For every $q>1$ there is a finite real constant $a_q$, and there is a finite real constant $A$, such that every Borel probability measure $\nu$ satisfies
\[
 a_q+\frac1q\operatorname{Ent}_{\mathrm{d}V}(\nu)
 \leq H_0(\nu)
 \leq A+n\operatorname{Ent}_{\mathrm{d}V}(\nu).
\]
The inequalities include the value $+\infty$. Thus $H_0(\nu)<\infty$ if and only if $\nu$ has finite entropy. A family of probabilities has uniformly bounded $H_0$ if and only if its entropies are uniformly bounded. Neither mass property is assumed.
\end{proposition}
\begin{proof}
In dimension one the closed formula \eqref{eq:normalization-closed-formula} gives $H_0=\log\int_X\omega+\operatorname{Ent}_{\mathrm{d}V}$, so the assertions follow. Suppose $n\geq2$.

For $G\in C(X)$, apply \cref{thm:lp-potential-estimate} to a bounded solution with density $\mathrm{e}^{B_0(G)+G}$. It gives a constant $c_q>0$, independent of $G$, such that
\[
 B_0(G)\geq\log c_q-\frac1q\log\int_X\mathrm{e}^{qG}\,\mathrm{d}V.
\]
For the other bound, fix a smooth Gauduchon metric $\gamma$, rescaled so that $\int_X\omega\wedge\gamma^{n-1}=1$, and write
\[
 k_\gamma=\left(\frac{\gamma^n}{\mathrm{d}V}\right)^{(n-1)/n},
 \qquad m=\min_X k_\gamma>0.
\]
For smooth $G$, the smooth equation, matrix arithmetic--geometric mean and Stokes give
\[
 \mathrm{e}^{B_0(G)/n}\int_X\mathrm{e}^{G/n}k_\gamma\,\mathrm{d}V\leq1.
\]
Uniform approximation and continuity of $B_0$ extend this inequality to continuous $G$. Consequently
\[
 B_0(G)\leq-n\log m-n\log\int_X\mathrm{e}^{G/n}\,\mathrm{d}V.
\]

If $\nu$ is not absolutely continuous, \cref{thm:normalization-dual-compactness} gives $H_0(\nu)=+\infty$, and both displayed bounds have their asserted extended-value meaning. For an absolutely continuous probability, the elementary variational calculation in the proof of \cref{prop:normalization-dual-entropy} gives
\[
 \sup_{F\in C(X)}\left(\int_XF\,\mathrm{d}\nu
                 -\log\int_X\mathrm{e}^F\,\mathrm{d}V\right)
       =\operatorname{Ent}_{\mathrm{d}V}(\nu),
\]
including infinite entropy. This calculation is measure-theoretic and does not require a closed background. Add $\int_XG\,\mathrm{d}\nu$ to the two bounds for $B_0(G)$ and take suprema. The substitutions $F=qG$ and $F=G/n$ give the claim with $a_q=\log c_q$ and $A=-n\log m$. Fixing, for example, $q=2$ proves both assertions about finite domains and uniformly bounded families.
\end{proof}

In particular the compact dual sublevels of \cref{thm:normalization-dual-compactness} have uniformly bounded entropy in this Hermitian zero-reference regime. This does not identify $H_0$ with entropy, bound the probability densities in a fixed $L^p$ space, or extend the assertion to an arbitrary singular reference.
